\chapter{Future Work and Open Problems}\label{ch:future-work}

The limitations weighed in \cref{ch:discussion} are not all equal: some are
permanent consequences of the central trade, while others are boundaries the
current system draws conservatively and could move with more work. This chapter is
concerned with the second kind. Each direction below is grounded in a limit the
system actually has, and each is described with an honest estimate of what it would
take and --- crucially --- what it would risk, because several of these extensions
push directly against the static-determinism property that the soundness gate and
the differential rig depend on. A direction that quietly sacrifices that property
would not be an extension of this work but a different system; where that danger
exists, it is named.

\section{Prescriptive scheduling}\label{sec:fw-prescriptive}

The schedule today is hand-written and the compiler is faithful to it; the
\texttt{check} assertions observe whether a bound is met but do nothing to meet it
(\cref{sec:lang-schedule}). The natural next step is to close that loop: turn a
checkable latency assertion into a solved latency budget, so that a schedule states
a goal --- a per-loop latency ceiling, a memory footprint, a worker count --- and
the compiler searches the schedule space for a decomposition that satisfies it.

This is the direction in which the compute-scheduling languages have invested most
heavily, through cost models and autotuners, and the same machinery is applicable
here. But it is genuinely hard, and it changes the contract: a solver needs a cost
model of each target, and a cost model is exactly the thing this work deliberately
does not have, because a faithful one is target-specific and a wrong one would
silently choose bad schedules. A credible version would likely begin narrowly ---
solving for one parameter, such as block size or pipeline depth, against a measured
rather than modelled cost on one target --- and treat full multi-objective schedule
synthesis as a long-term goal rather than a near-term feature. The value of the
present architecture for this work is that the search space is already explicit and
the soundness gate already rejects unsound candidates, so a solver could explore
schedules knowing that every accepted one is sound by construction.

\section{Floating-point determinism}\label{sec:fw-float}

The bit-identity guarantee once stopped at floating-point summation, because IEEE-754
addition is not associative and a distributed float-sum reduced in a varying order
produces varying bytes (\cref{sec:disc-shortcomings}). That boundary has now been
crossed for the cross-backend case, but only partway, and the remaining distance is
worth stating precisely.

What is implemented is the \emph{fixed-order} route: an opt-in \texttt{combine=fsum}
under which the host folds the per-worker partials in a fixed worker-id-sorted order
that is identical on every backend, and the reference oracle reproduces that exact
fold. The reduction lowering, the kernel contract, and the reference implementation
therefore move in lockstep --- the subtlety anticipated below, made concrete --- and a
distributed float sum now carries the same cross-backend byte-identity the integer
combines enjoy, demonstrated across all seven CPU backends by the \texttt{28-bin-fsum}
example. The plain \texttt{sum} combine stays rejected, so the guarantee is opt-in and
the default stays strict.

What remains open is the stronger, \emph{order-independent} route: a reproducible
reduction that computes a fixed result regardless of summation order, for instance by
accumulating into a wide fixed-point superaccumulator --- an integer accumulator wide
enough to represent any partial sum exactly --- that absorbs each addend without
rounding, or by pre-rounding addends to a common
granularity~\cite{demmel2013reproducible}. The fixed-order fold delivered here is
reproducible across backends for a \emph{given} schedule, but it is not the naive
left-to-right IEEE-754 sum, and a different worker count folds the partials
differently --- so it does not yet give ``same bytes across schedules.'' An
order-independent scheme would close that gap, making every decomposition agree
bit-for-bit on a single reproducible result independent of worker count (the result
being the scheme's own --- the exact sum for a superaccumulator, the pre-rounded sum
for pre-rounding --- not necessarily the naive single-pass IEEE sum). The subtlety that
made even the fixed-order step less self-contained than it first appears --- that the
immutable oracle, the rig's trust anchor, must adopt the same scheme so cross-backend
agreement is also agreement with the oracle --- applies with full force to the
order-independent route, and is the reason it is the higher-value but heavier
remaining direction: it lifts the schedule-invariance limit that still excludes part
of a whole numerical domain, but it is not the trivial drop-in the opt-in framing
alone would suggest.

\section{Beyond the affine, static class}\label{sec:fw-affine}

The hardest and most consequential limits are the ones that define the algorithm
class itself, and they come in several grades of difficulty.

\paragraph{Distributed gather and scatter.} Data-dependent indexing is supported
today only conservatively: a gather broadcasts the whole indexed array and a
scatter replicates and recombines its target, and a scatter onto an affine-
partitioned target is rejected (\cref{sec:arch-transfer}). A precise distributed
gather/scatter --- one that moves only the elements a worker actually touches ---
would require a communication-synthesis pass that reasons about the index array's
contents, which are not known at compile time, and so would need either a runtime
exchange phase or a two-pass inspector/executor scheme --- a first pass that
inspects the index array to plan the communication and a second that executes that
plan --- of the kind sparse libraries use~\cite{saltz1991runtime}. This is well-trodden ground in distributed sparse computing, but importing it
here means admitting a runtime-determined communication pattern, which sits in
tension with the static firing order. The difficulty is sharper than the
bounded-iteration case treated below in this section. Bounded data-dependent
\emph{iteration} (the capped \texttt{until} loop exercised by
\texttt{21-jacobi-converge} and \texttt{29-jacobi-cap-hit}) works because its only
runtime unknown is a scalar trip count: a compile-time cap fixes the worst-case
firing order and any shorter run replays as a sub-trace of it. A precise
gather/scatter's runtime unknown is not a scalar but an entire address map --- which
element moves to which worker --- whereas the soundness gate replays a communication
structure fixed at compile time. The only compile-time-fixed structure that covers
every possible index pattern is the conservative whole-array broadcast the backends
already emit; that broadcast is, in effect, the trivial statically sound envelope,
merely an untightened one. Tightening it is therefore not blocked by bounding the
exchange's \emph{size} --- that bound is the part the iteration work already
demonstrates --- but by deciding \emph{which} elements move while the gate still sees
a fixed structure. Closing the gap needs either a gate that admits a bounded family
of runtime communication patterns or a static envelope provably sound for any pattern
yet tighter than whole-array; both are verification advances the present soundness
argument does not reach, which is why the conservative broadcast remains the only
sound option today.

\paragraph{Data-dependent control flow.} Convergence-driven termination --- a loop
that runs ``until converged'' --- is the single most-requested capability the model
still lacks in general, and lifting it \emph{fully} is the deepest change, because an
unbounded data-dependent trip count means the firing order is no longer statically
determined, which is the property everything else rests on. The bounded middle path,
however, is no longer hypothetical. Bounded data-dependent iteration --- a loop with a
compile-time maximum trip count and a runtime early exit, so the worst-case firing
order remains statically analysable even though the actual run may stop early --- is
\emph{realized on the single-worker path}. Two corpus examples
(\texttt{21-jacobi-converge} and \texttt{29-jacobi-cap-hit}, the optional
\texttt{until} clause of \cref{app:grammar}) compile and run a Jacobi convergence
loop value-correctly: one exits early, and one runs the full cap so the worst-case
bound the soundness gate analyses is itself an exercised, byte-identical execution
path rather than an assertion. This holds because the gate replays the full capped
unroll and any early exit is a sub-trace of it (soundness preserved by the cap), and
because the exact-integer halt predicate is a deterministic, backend-agnostic
function of the data, so every backend stops at the same generation (the differential
test preserved). Two residuals remain. First, the break is single-worker only: on the
multi-worker path the convergence scalar is a cross-worker reduction, so the early
exit must become a collective all-reduce-and-broadcast at a barrier --- every worker
branching on the identical global value and exiting at the same generation by
construction --- which keeps the firing order static but is a genuinely new collective
pattern not yet built; multi-worker convergence cells are recorded as skips meanwhile.
Second, byte-identity is at risk under a \emph{floating-point} convergence predicate,
whose comparison value depends on accumulation order, tying this direction back to the
float-determinism limit above. Even so, the realized form already admits a useful
subset of single-worker iterative solvers --- those with a known iteration cap and an
exact termination test --- and the residual is a bounded, well-scoped lift rather than
the open-ended problem unbounded data-dependent control flow poses.

\paragraph{Sparse iterative solvers.} A full sparse solver needs both of the
above --- data-dependent indexing and convergence-driven termination --- which is
why it sits entirely outside the present model. It is listed here not as a
near-term target but as the natural destination if the two preceding directions
both succeed, and as an honest marker of how far the current class is from
general numerical computing.

\paragraph{Precise inference, and the remaining halves of the region saving.} The
groundwork here is now partly built, which sharpens what is left to do. The backends
already transmit only the inferred region on a cross-worker edge whose region is a
single contiguous span, cutting wire volume measurably (\cref{sec:res-quant-comm}); two
distinct extensions remain, one on the inference side and one on the backend side, and
they are independent. On the inference side, the transfer-inference pass still degrades
to whole-array broadcast whenever an index is not an affine, contiguous function of the
partition variables. A polyhedral analysis, using an established integer-set library,
could compute precise transferred regions for a broader class of affine and
piecewise-affine index expressions than the current prefix-based inference handles,
without changing the language --- and because the wire now follows the inferred region,
a sharper inference would directly narrow more edges rather than merely improving
receiver-side correctness. On the backend side, two cases the wire flip deliberately
left whole-array remain: a \emph{strided} or non-prefix region (a column band of a
row-major matrix, an irregular halo) needs a banded pack-and-scatter the backends do not
yet emit, and the embedded tier needs a banded receive to match a banded send. And the
\emph{footprint} half is wholly untouched: every worker still allocates the array
full-shaped and pastes its received span into place, so per-worker memory is unchanged;
reducing the allocation to the owned band needs index rebasing at kernel-call sites or a
banded-view abstraction, because downstream reads index with global coordinates
(\cref{sec:disc-shortcomings}). All of these are internal precision improvements that
risk nothing in the guarantees --- the firing order, the soundness gate, and the
byte-level results are unchanged, only the volume moved and the memory held --- which
keeps the direction comparatively safe, but it is a several-part one, of which only the
contiguous-span wire flip is so far done.

\section{New target tiers}\label{sec:fw-tiers}

The three tiers --- commodity CPU, message-passing cluster, embedded
microcontroller --- cover a wide span but stop short of massively parallel
accelerators. A GPU or neural-processing-unit tier is a plausible fourth, and the
event contract is in principle agnostic enough to drive one. But the fit is not
free: a GPU's execution model is bulk-synchronous (alternating compute phases with global
synchronisations across many parallel threads) over thousands of lanes, with a
memory hierarchy the current contract does not model, and the \texttt{Alloc} and
\texttt{Free} events (\cref{sec:arch-petri}), currently reserved in the event
contract but emitted by no backend, would have to become load-bearing to place data
across that hierarchy. The honest assessment is that a
GPU tier is a research project in its own right, not a matter of writing one more
shim, and that claiming the architecture ``extends to GPUs'' without having done it
would be exactly the kind of overclaim this document has tried to avoid. It is
listed as an open direction, not a promised one.

The cheaper and more certain tier work is further breadth within the existing
embedded tier. A second microcontroller family has already been added --- an
nRF52840 (Cortex-M4F) shim whose UARTE EasyDMA transmit is a genuinely different
register interface from the STM32 USART --- byte-exact on the same examples
(\cref{ch:results}), which both demonstrates the hardware-abstraction trait is a
real portability seam and calibrates the cost: that second family reused the generic
backend and the test harness unchanged, swapping only the concrete shim, the memory
map, and one linker flag, because it shared the Cortex-M target the toolchain
already builds and had a ready simulator model. Additional families --- each a
separate shim against the committed trait --- would broaden the evidence further.
This is bounded, well-specified engineering rather than research; cheap when a family
shares the target and has a simulator model, more expensive when it needs a new
target triple or a model that does not yet exist. It remains the most direct way to
widen the part of the claim that \cref{sec:disc-shortcomings} identifies as the
narrowest.

\section{A quantitative performance and scaling study}\label{sec:fw-quant}

This document confined its main results to correctness and coverage, because the
contribution it defends is the separation and its falsifiability, not the speed of
the emitted code (\cref{ch:results}). It does report a set of \emph{static}
engineering characteristics --- per-backend code footprint, generated-code size, and
how the soundness gate's analysis net scales with problem size
(\cref{sec:res-quant}) --- so the cost of the genericity is at least described. Two
quantitative directions remain genuinely open, and it is worth stating each
precisely.

The first is a \emph{runtime} performance and scaling study: how the emitted
programs' wall-clock time varies with worker count, schedule, and backend on real
hardware --- and, for the cluster tier, on an actual cluster rather than a local
launcher. This is the measurement that would turn the portability claim from ``the
same values everywhere'' into ``the same values, and here is what each decomposition
costs,'' and it is the input the prescriptive-scheduling direction of
\cref{sec:fw-prescriptive} would need. It is deliberately deferred here because a
credible runtime study demands a benchmarking methodology --- worst-case inputs,
controlled hardware, warm-up and variance discipline --- that is a project in itself,
and because no runtime number bears on the correctness claims already established: a
slow but byte-identical program still falsifies nothing.

A scoping observation explains why this was, until recently, more than a matter of
securing time on a cluster, and it once coupled the runtime study tightly to the
gate-scaling direction below. At the corpus's problem sizes the emitted programs'
wall-clock is dominated by fixed per-invocation overhead rather than by computation ---
a single-worker matrix multiply at the corpus dimension finishes in about a
millisecond, of which the few thousand integer multiply--adds account for only
microseconds; the remainder is process spawn, input loading, and thread setup, not
arithmetic --- so a meaningful scaling signal emerges only at substantially larger
problem sizes. The gate is not separable from the build: there is no option to emit a
decomposed program without analysing it, so the cost of compiling the decompositions
one would want to time used to bound what could be timed at all. That coupling has now
loosened for the dominant case. The gate proves the single-shot matched-pair
communicating subclass --- which is the shape of most of the corpus's distributed
schedules --- symbolically, so building the four-worker distributed matrix multiply at
$N=512$ now costs about 7.2\,megabytes of compiler memory rather than the hundred-plus
gigabytes its expanded net would project (\cref{sec:res-quant-net}). One can therefore
now compile exactly those large distributed decompositions, which removes the gate as
the binding obstacle to timing them. What remains coupled is narrower: the loop-carried
and pipelined distributed shapes still expand, so a runtime study that needed to scale
\emph{those} specific schedules would still wait on the gate work below. A credible
runtime study therefore now waits primarily on a representative cluster and a
disciplined methodology rather than on the compiler --- except for the residual
loop-carried and pipelined shapes, where the two open problems remain coupled.

The second direction is the gate-scaling limitation that \cref{sec:res-quant-net}
measured, and it has been \emph{substantially} closed, in two steps. Because the
soundness gate replayed a net expanded over the iteration space --- roughly one place and
one transition per kernel firing --- its cost was linear in the number of firings, cheap
for the corpus but unable to scale to a problem with billions of firings. First, for the
\emph{buffer-free} subclass (a single worker, or any schedule with no cross-worker
transfer) the gate now reasons about the \emph{looped} structure directly: it proves
soundness from the rolled ACFG in time independent of the iteration count, never building
the expansion. Second, and the harder step, the symbolic proof now reaches across a
cross-worker boundary, to the \emph{single-shot matched-pair} communicating subclass: a
distributed schedule whose every transfer is one \texttt{Push} matched by one
\texttt{Wait} at loop depth zero, so the shared buffer place's peak occupancy is provably
one and at most its capacity. This is the dominant distributed shape --- the corpus's
distributed reductions, stencils, matrix multiply, sparse mat--vec, histogram, and
scatter all fall in it --- and proving it symbolically, against a firing-order invariant
pinned by property tests and cross-checked against the expanded replay on every gated
cell, is what keeps the $N=512$ distributed matmul compile flat. What remains open is the
\emph{loop-carried} and \emph{pipelined} communicating shapes --- a transfer whose buffer
occupancy depends on cross-iteration drainage, or one that deliberately keeps several
frames in flight --- whose boundedness still depends on a firing-order interleaving the
symbolic analysis does not yet cover. A steady-state occupancy argument over their looped
\texttt{Push}/\texttt{Wait} structure is the concrete compiler work still outstanding. It
is deliberately not approximated: an imprecise buffer analysis could accept a net the
expanded replay would reject, trading soundness for scaling, which the gate's whole
purpose forbids --- so anything outside the proven subclasses is routed loudly to the
sound expanded replay until the symbolic argument is extended to it too.

Taken together, these directions share a common discipline: each is judged not only
by what it would add but by what it would risk to the determinism the system is
built on. Two of them sit off that risk axis entirely --- prescriptive scheduling
and the quantitative study change no guarantee, the first being an optimisation
question and the second an empirical one. Of the rest, the order-independent float
reduction and the polyhedral precision improvement add capability without touching
the guarantees, even though the polyhedral work is grouped, by topic, inside the
affine-class section; bounded data-dependent iteration and a precise distributed
scatter push against the static firing order and must be done carefully if at all;
and unbounded data-dependent control flow and a GPU tier would, if done naively,
dismantle the very properties that make the present system worth having. The
interesting future work is therefore not simply to lift the limits, but to lift them
while keeping the soundness gate sound and the differential test meaningful --- which
is the same constraint that shaped the original design.
